\documentclass[a4paper,9pt]{extarticle}

%=========================================================
% PACKAGES
%=========================================================

\usepackage[
a4paper,
top=1.8cm,
bottom=1.5cm,
footskip=7mm,
left=1cm,
right=1cm,
headheight=14pt,
headsep=5mm
]{geometry}
\usepackage[T1]{fontenc}
\usepackage[utf8]{inputenc}
\usepackage[french]{babel}

\usepackage[sfdefault]{FiraSans}
\renewcommand{\familydefault}{\sfdefault}

\usepackage{multicol}
\setlength{\columnsep}{0.7cm}

\usepackage{amsmath,amssymb}

\usepackage{titlesec}
\usepackage{enumitem}

\usepackage{tcolorbox}
\tcbuselibrary{skins,breakable}

\usepackage{xcolor}

\usepackage{graphicx}
\usepackage{float}

\usepackage{tikz}

\usepackage{fancyhdr}
\usepackage[hidelinks]{hyperref}


%=========================================================
% PAGE SETTINGS
%=========================================================

\setlength{\parindent}{0pt}
\setlength{\parskip}{1.5pt}

\setlength{\abovedisplayskip}{4pt}
\setlength{\belowdisplayskip}{4pt}

\setlist[itemize]{
leftmargin=*,
itemsep=1pt,
topsep=1pt
}


%=========================================================
% COLORS
%=========================================================

\definecolor{Blue}{HTML}{E3F0FF}
\definecolor{Green}{HTML}{E3F7EA}
\definecolor{Pink}{HTML}{FFE5ED}
\definecolor{Orange}{HTML}{FFF0DC}
\definecolor{Purple}{HTML}{EEE7FF}
\definecolor{Yellow}{HTML}{FFFADB}


\definecolor{BlueTitle}{HTML}{285B8F}
\definecolor{GreenTitle}{HTML}{287052}
\definecolor{PinkTitle}{HTML}{A84063}
\definecolor{OrangeTitle}{HTML}{A96500}
\definecolor{PurpleTitle}{HTML}{6245A5}


%=========================================================
% HEADER
%=========================================================

\pagestyle{fancy}

\fancyhf{}

\fancyhead[L]{\small\textbf{Responsible Software}}
\fancyhead[C]{\small Safety 1: Harms and Threats}
\fancyhead[R]{\small Week 04 --- 28/09-04/10/2026}

\fancyfoot[C]{\small\href{https://github.com/AlphaXT38G/notes_ba5}{github.com/AlphaXT38G/notes\_ba5}}
\fancyfoot[R]{\small \thepage}

\renewcommand{\headrulewidth}{0.4pt}



%=========================================================
% TITLES
%=========================================================

\titleformat{\section}
{\large\bfseries\color{BlueTitle}}
{}{0pt}{}

\titlespacing{\section}
{0pt}{6pt}{3pt}


% separator

\newcommand{\separator}{
\vspace{2pt}
\hrule
\vspace{3pt}
}


%=========================================================
% BOXES
%=========================================================

% Boxes are kept whole (unbreakable): inside multicols, a box split by tcolorbox
% leaves an empty column behind it.

% DEFINITION - blue
\newtcolorbox{definition}{
enhanced,
unbreakable,
colback=Blue,
colframe=BlueTitle,
boxrule=0.7pt,
arc=0mm,
title=\textbf{Definition},
fonttitle=\bfseries
}

% IMPORTANT - pink
\newtcolorbox{important}{
enhanced,
unbreakable,
colback=Pink,
colframe=PinkTitle,
boxrule=0.7pt,
arc=0mm,
title=\textbf{Key idea},
fonttitle=\bfseries
}

% FORMULAS - green
\newtcolorbox{formula}{
enhanced,
unbreakable,
colback=Green,
colframe=GreenTitle,
boxrule=0.7pt,
arc=0mm,
title=\textbf{Key formulas},
fonttitle=\bfseries
}

% WARNING - orange
\newtcolorbox{warning}{
enhanced,
unbreakable,
colback=Orange,
colframe=OrangeTitle,
boxrule=0.7pt,
arc=0mm,
title=\textbf{Warning},
fonttitle=\bfseries
}

% EXAMPLE - purple
\newtcolorbox{example}{
enhanced,
unbreakable,
colback=Purple,
colframe=PurpleTitle,
boxrule=0.7pt,
arc=0mm,
title=\textbf{Example},
fonttitle=\bfseries
}

% SUMMARY - orange
\newtcolorbox{summary}{
enhanced,
unbreakable,
colback=Yellow,
colframe=OrangeTitle,
boxrule=0.7pt,
arc=0mm,
title=\textbf{Summary},
fonttitle=\bfseries
}



% GOAT - brown: the idea told as simply as possible
\definecolor{Goat}{HTML}{F7EEDF}
\definecolor{GoatTitle}{HTML}{8A5A1F}
\newtcolorbox{goat}{
enhanced,
unbreakable,
colback=Goat,
colframe=GoatTitle,
boxrule=0.7pt,
arc=0mm,
title=\textbf{Explained to a goat},
fonttitle=\bfseries
}



%=========================================================
% FIGURE BOX
%=========================================================

\newtcolorbox{figurebox}{
enhanced,
colback=white,
colframe=gray,
boxrule=0.5pt,
arc=0mm,
center
}



%=========================================================
% DOCUMENT
%=========================================================

\begin{document}

%\footnotesize   --> To reduce the size of the text

\thispagestyle{fancy}



\begin{multicols}{2}


%=========================================================
\section{Safety and Security}


\begin{definition}

\textbf{SECURITY vs. SAFETY} (Leveson, 1986)

\begin{itemize}

\item \textbf{Security}: protect the \textbf{system} from the \textbf{environment};
\textbf{threats}, \textbf{intentional}. Methods: vulnerability scans, ethical hacking, audits,
code review.

\item \textbf{Safety}: protect the \textbf{environment} from the \textbf{system};
\textbf{hazards}, \textbf{accidental}. Methods: testing, failure mode analysis, redundancy.

\end{itemize}

Both are about \textbf{risks} (probability of something bad).

\end{definition}


\begin{example}

\begin{itemize}

\item Car misreads 10\% of signs: \textbf{safety hazard}.

\item Stickers make a stop sign read as a speed limit: \textbf{security threat} +
\textbf{safety hazard}.

\end{itemize}

\end{example}


\textbf{Socio-technical system}: social and technical systems affect each other (social media,
autonomous cars): hidden loops, emergence, non-determinism.


\begin{goat}

Security: stop the wolf breaking the fence (on purpose). Safety: stop the fence falling on the
goats (by accident).

\end{goat}



%=========================================================
\section{Four Scenarios of Harm}


\begin{definition}

\textbf{WHEN SOFTWARE CAUSES HARM} (Parasuraman \& Riley)

\begin{itemize}

\item \textbf{Malfunction}: does not work as specified, unintentional bug (Therac-25,
CrowdStrike 2024).

\item \textbf{Abuse}: malicious use by bad actors (ransomware on hospitals).

\item \textbf{Misuse}: unintended / unsupported use, careless (Samsung code in ChatGPT).

\item \textbf{Intended use}: used as planned and still harms (infinite scroll, elections);
the hardest to anticipate.

\end{itemize}

\end{definition}


\begin{warning}

Misuse vs. abuse: the difference is \textbf{intent}. Viral travel trends ruining ecosystems:
\textbf{intended use}.

\end{warning}


\begin{goat}

It breaks, a villain uses it, someone uses it wrong, or it works perfectly and still hurts.

\end{goat}



%=========================================================
\section{User-Generated Content}


\begin{definition}

\textbf{UGC}: content (text, images, video) posted publicly by users on platforms.

\textbf{Harms} (Banko 2020): hate \& harassment, self-inflicted harm, ideological harm,
exploitation. \textbf{Toxic content}: harmful or offensive material.

\end{definition}


\begin{definition}

\textbf{HATE SPEECH} (UN 2019): discriminatory or pejorative + about identity factors
(religion, ethnicity, gender...) + any form of expression.

\end{definition}


\begin{example}

\textbf{Facebook and the Rohingya} (Myanmar 2017): hate and military disinformation, weak and
slow moderation.

\end{example}


\begin{important}

\textbf{Free speech} (UDHR art. 19) is not absolute. Filtering errors:

\begin{itemize}

\item \textbf{FN}: toxic content escapes (people harmed);

\item \textbf{FP}: harmless content censored (free speech harmed).

\end{itemize}

Who decides? States, platforms, people.

\end{important}


\begin{goat}

Delete too little and people get hurt; delete too much and innocent voices are silenced.

\end{goat}



%=========================================================
\section{Content Moderation}


\begin{definition}

\textbf{CONTENT MODERATION} (Grimmelmann 2015): governance mechanisms structuring participation
to facilitate cooperation and prevent abuse = \textbf{identification} + \textbf{action}.

\end{definition}


\begin{formula}

\textbf{When}: pre-publication; post-publication (systematic, sampling, \textbf{reactive} on
reports).

\textbf{How}:

\begin{itemize}

\item \textbf{Classification}: keyword screening, sentiment analysis, ML classifiers, computer
vision.

\item \textbf{Matching}: hashing vs. database; \textbf{perceptual hashing} robust to small
changes (Christchurch video: 80\% of re-uploads blocked).

\end{itemize}

\textbf{Actions}: \textbf{exclude} (ban, timeout) $>$ \textbf{organize} (delete, annotate,
reduce visibility) $>$ \textbf{set norms} (warn).

\end{formula}


\begin{warning}

Automation misses context and nuance; users evade filters $\Rightarrow$ \textbf{hybrid}. Human
moderators face severe mental health risks: a \textbf{dilemma} (some exposed to protect
others).

\end{warning}


\begin{goat}

Find bad posts (by type or by matching known ones), then ban, delete, hide or warn. Machines
miss jokes; the humans who help get hurt.

\end{goat}



%=========================================================
\section{The Confusion Matrix}


\begin{formula}

Rows = \textbf{actual}, columns = \textbf{predicted}:

\begin{center}
\begin{tabular}{l|c|c}
 & pred. + & pred. $-$ \\ \hline
actual + & TP & FN \\ \hline
actual $-$ & FP & TN \\
\end{tabular}
\end{center}

\[
\text{Acc} = \frac{TP+TN}{\text{all}}
\]
\[
\text{FNR} = \frac{FN}{TP+FN} \qquad \text{FPR} = \frac{FP}{TN+FP}
\]

\end{formula}


\begin{example}

\textbf{Fissures} (positive = fissure): TP 60, FN 15, FP 20, TN 100.

FNR $= 15/75 = 20\%$. For safety, watch \textbf{FN} (missed fissures).

\end{example}


\begin{warning}

Check what \textbf{positive} means. COVID test: FN dangerous. Quality control with positive =
``good'': FP dangerous.

\end{warning}


\begin{goat}

False alarm (FP) or missed wolf (FN): which is worse depends on the story.

\end{goat}



%=========================================================
\section{Strategy: Bad Actors}


\begin{definition}

\textbf{BAD ACTOR}: person or group with \textbf{malicious intent} exploiting a system; often
external, also insiders.

\end{definition}


\begin{formula}

\textbf{5 motivations} (they overlap):

\begin{itemize}

\item \textbf{Money}: scammers, hackers, competitors (ransomware, data theft).

\item \textbf{Politics}: nation states, agencies (disinformation, infrastructure attacks).

\item \textbf{Entertainment}: trolls, thrill seekers.

\item \textbf{Ideas}: hacktivists, extremists, conspiracy theorists (DoS, leaks, defacement).

\item \textbf{Self-interest}: revenge, obsession, insiders (harassment, tracking, grade
changing).

\end{itemize}

\end{formula}


\begin{warning}

Graded: \textbf{intent} explicit + \textbf{motivation} explicit + chain \textbf{action
$\rightarrow$ impact on stakeholders}, specific and diverse.

\end{warning}


\begin{goat}

Imagine the villains, why they come, what they do, who gets hurt.

\end{goat}


\end{multicols}


\end{document}
